\section{Introduction}
\label{sec:introduction}

Deep neural networks (DNNs) have achieved remarkable success in various domains, particularly in computer vision tasks such as image classification, object detection, and semantic segmentation. However, their vulnerability to adversarial examples---carefully crafted inputs designed to induce misclassification---remains a critical security concern~\cite{szegedy2014intriguing, goodfellow2015explaining}. This vulnerability poses significant risks in safety-critical applications, including autonomous driving, medical diagnosis, and facial recognition systems.

Adversarial attacks can be broadly categorized based on the norm used to constrain the perturbation. While $\ell_2$ and $\ell_\infty$ attacks have been extensively studied~\cite{madry2018towards, carlini2017towards}, $\ell_0$-norm attacks---which minimize the \emph{number} of perturbed pixels---have received comparatively less attention due to the non-differentiable and combinatorial nature of the $\ell_0$-norm. Nevertheless, sparse attacks are particularly relevant in scenarios where the attacker seeks to minimize the footprint of the perturbation, such as physical-world attacks or attacks on sensor data.

The recently proposed $\sigma$-zero algorithm~\cite{cina2025sigma} represents a significant advance in $\ell_0$-norm adversarial attacks. By employing a smooth approximation of the $\ell_0$-norm combined with dynamic thresholding, $\sigma$-zero achieves state-of-the-art attack success rates while maintaining sparse perturbations. The key insight is to replace the non-differentiable $\ell_0$-norm with a continuous surrogate:
\begin{equation}
    \norm{\deltaVec}_0 \approx \sum_{i} \frac{\delta_i^2}{\delta_i^2 + \sigma},
\end{equation}
where $\sigma > 0$ controls the approximation fidelity.

\subsection{Motivation and Problem Statement}

Despite its effectiveness in minimizing the $\ell_0$-norm, we identify a critical limitation of the original $\sigma$-zero algorithm: \emph{the generated adversarial examples are often visually detectable}. This occurs because the algorithm optimizes solely for sparsity (number of perturbed pixels) without constraining the \emph{magnitude} or \emph{spatial distribution} of the perturbations. Specifically:

\begin{enumerate}
    \item \textbf{Unbounded per-pixel changes:} Even when only a few pixels are modified, those pixels may undergo drastic changes (e.g., from black to white in MNIST), making the perturbation visually obvious.
    \item \textbf{No structural awareness:} The algorithm does not consider the spatial structure of the image, potentially disrupting important visual features.
    \item \textbf{No spatial smoothness:} Perturbations may be scattered randomly across the image, creating visible artifacts.
\end{enumerate}

These limitations are illustrated in Figure~\ref{fig:visual_comparison}, which compares adversarial examples generated by the original $\sigma$-zero and our proposed method.

\subsection{Contributions}

In this paper, we propose \textbf{Perceptually-Aware $\sigma$-zero}, an enhanced algorithm that addresses the visual quality limitations of the original $\sigma$-zero. Our contributions are:

\begin{enumerate}
    \item \textbf{Visual perception constraints:} We introduce three complementary regularizers---$\ell_\infty$-norm penalty, SSIM loss, and TV regularization---into the $\sigma$-zero optimization framework.
    \item \textbf{Adaptive weighting scheme:} We propose a dynamic weighting strategy that balances attack aggressiveness and visual quality throughout optimization.
    \item \textbf{Multi-scale SSIM:} We extend the SSIM loss to multiple scales for more comprehensive structural preservation.
    \item \textbf{Comprehensive evaluation:} We conduct extensive experiments on MNIST, demonstrating significant improvements in visual quality (SSIM +9.5\%) while maintaining 100\% attack success rate and achieving even sparser perturbations (median $\ell_0$: 23 $\to$ 18).
    \item \textbf{Ablation and sensitivity analysis:} We provide detailed analysis of each component's contribution and hyperparameter sensitivity.
\end{enumerate}
